\chapter*{Resumo}
\markboth{Resumo}{Resumo}

O sistema pericial de condução decide, a cada 100\,ms, a velocidade e a faixa
de um veículo simulado. A simulação (física, sensores e mundo) está em Java. A
política de condução vive num motor de regras, e existem dois motores
intermutáveis: Drools e SWI-Prolog. Ambos implementam a interface
\cod{ExpertSystemBackend}, pelo que o \cod{SimulationEngine} não sabe qual está
ativo. O motor escolhe-se no arranque (\cod{expert.backend}) ou durante a
execução (\cod{POST /api/simulation/backend/prolog}).

A base de conhecimento Prolog replica todas as regras de decisão do Drools com
os mesmos nomes e com a saliência do Drools como precedência. Tem 120 regras:
75 com o nome exato de uma regra Drools, 43 variantes por resultado de 16 regras
Drools (por exemplo \cod{'Hold at red traffic light (emergency)'}) e 2 regras
intermédias só do Prolog. Este documento descreve o motor Prolog, a sua ligação
ao resto do sistema, o metaconhecimento dinâmico e as explicações.

Este relatório descreve um sistema pericial de apoio à condução que recebe
observações de sensores, mantém estado temporal e produz decisões explicáveis
sobre velocidade e mudança de faixa. A componente central analisada é o
\emph{backend} SWI-Prolog. O projeto não utiliza apenas o mecanismo de resolução
nativo para responder a consultas: implementa um motor próprio de regras de
produção com memória de trabalho, agenda, precedência, refração, encadeamento
para a frente e reavaliação seletiva através de um índice de dependências.

A base de conhecimento contém 120 regras declarativas que cobrem travagem de
emergência, peões, semáforos, sinais STOP, limites de velocidade, seguimento de
veículos, ultrapassagem, obstáculos, buracos, barreiras e condições
meteorológicas com lógica difusa. O motor combina estas regras com 38 tipos de
estado dinâmico. Cada ciclo da simulação envia ao processo Prolog o estado do
veículo e os eventos sensoriais em JSON. O Prolog devolve propostas, factos
derivados, verificações de segurança lateral e métricas do motor; a camada Java
arbitra as propostas em dois canais independentes, longitudinal e lateral.

A explicabilidade resulta do registo da proveniência de cada facto e de cada
proposta. O sistema responde a perguntas do tipo como foi obtida a decisão, por
que motivo venceu uma regra e por que motivo uma alternativa não foi aplicada. A
validação inclui testes unitários, testes de sequências temporais e testes de
equivalência entre Prolog e Drools. Na execução verificada para este relatório,
o \emph{build} Maven executou 598 testes sem falhas (2 ignorados).

\vspace{0.8cm}
\noindent\textbf{Palavras-chave:} sistemas periciais; Prolog; motor de
inferência; regras de produção; metaconhecimento; explicabilidade; apoio à
condução.

\clearpage
